\chapter*{Resumen}
\addcontentsline{toc}{chapter}{Resumen}

Esta memoria estudia el pronóstico de series de tiempo no equiespaciadas mediante arquitecturas tipo \Transformer\ que reciben directamente los tiempos de observación, sin remuestrear las señales sobre una grilla fija. La investigación parte de \CTTransformer, un encoder con representación temporal continua y consultas futuras, y desarrolla dos variantes orientadas explícitamente al tiempo físico: \QueryCross, que codifica sólo la historia y resuelve cada tiempo objetivo mediante atención cruzada independiente con sesgos de lag relativo; y \ContinuousBasis, que representa el pronóstico como una función continua del horizonte mediante bases de tendencia, funciones radiales y Fourier. Ambas admiten una transición de estado continuo estable (\CTSSM) y cabezas puntuales o gaussianas heteroscedásticas.

Se distinguen dos evaluaciones que no son directamente comparables. El benchmark legado sobre cinco series reales usa historias y horizontes definidos en cantidad de eventos. Sus resultados muestran que las arquitecturas propuestas pueden entrenarse sobre timestamps irregulares y entregan evidencia descriptiva de competitividad, pero su ablación temporal no identifica por sí sola el aporte del tiempo físico: el orden de las consultas, la densidad de eventos y la persistencia pueden resolver parte de la tarea. Además, las semillas de un mismo dataset no se consideran réplicas científicas independientes; el dataset es la unidad experimental.

Para abordar estas limitaciones se especifica una evaluación física sintética congelada y versionada sobre procesos univariados y multivariados. El protocolo fija historias por duración física, obtiene los objetivos desde la misma trayectoria latente limpia almacenada por separado de las observaciones ruidosas, aleatoriza el orden de las consultas, conserva timestamps absolutos en doble precisión hasta el recentrado local y evalúa por horizonte, densidad, hueco máximo y edad de la última observación. Así, los targets son independientes del patrón de observación y del ruido, no de la realización latente. Incluye controles de identificabilidad, ablaciones temporales emparejadas, persistencia y métricas probabilísticas de calibración. Como los procesos y diagnósticos fueron inspeccionados durante el desarrollo, su alcance es descriptivo y no equivale a una prueba prospectiva preregistrada. El código genera automáticamente tablas LaTeX una vez completado el benchmark. Hasta entonces, las ejecuciones de humo sólo validan contratos de software y no establecen un ranking neuronal.

\noindent\textbf{Palabras clave:} series de tiempo irregulares, Transformers, tiempo continuo, atención cruzada, modelos de estado, pronóstico probabilístico.
